\documentclass[a4paper,10pt]{scrartcl}

\usepackage[utf8]{inputenc}
\usepackage{tikz}
\usepackage{amsmath}
\usepackage{hyperref}
\usepackage{import}
\usepackage{xcolor}

\usepackage{subfig}
\usepackage{graphicx}
\usepackage{pgfplots}
\pgfplotsset{compat=1.18}
\usepackage{pgfplotstable}
\tikzset{>=stealth}
\usetikzlibrary{patterns}
\usetikzlibrary{pgfplots.statistics}
\usepgfplotslibrary{external}
\tikzexternalize
\tikzsetexternalprefix{plots/}
%\tikzset{external/force remake}
%%%%%%%%%%%%%%%%%%%%%%% PGFPLOTS DEFAULTS
\pgfplotstableset{col sep=comma}
\pgfplotstableset{header=false}

%%%%%%%%%%%%%%%%%%%%%%% COLORS
\definecolor{KKLightGreen}{RGB}{35,248,222}
\definecolor{KKGreen}{RGB}{30,140,153}
\definecolor{KKDarkGreen}{RGB}{14,130,130}

\definecolor{KKLightPurple}{RGB}{201,49,161}
\definecolor{KKPurple}{RGB}{99,41,199}
\definecolor{KKDarkPurple}{RGB}{25,23,123}

%%%%%%%%%%%%%%%%%%%%%% NOTES in red
\newcommand\gist[1]{\textcolor{red}{#1}}

\title{Influence of the Heterogenious Memory Management System on Kokkos Semantics }

\author{Jakob Bludau}


\begin{document}
\maketitle

\section{Introduction}
\gist{
  Prior Knowledge:
\begin{itemize}
\item Knowledge about Heterogenious computing being a thing
\end{itemize}
  What is introduced in this chapter:
\begin{itemize}
\item Performance benefits from heterogeneous computing
\item most HPC machines heterogeneous
\item Heterogenious computing uses multiple, disjunct compute resources
\item Heterogenious computing uses (potentially) different hardware locations for memory
\item For successful execution the compute resource has to have access to the memory
\end{itemize}
}

\section{Memory Management in Heterogenious Architectures}
\gist{
  Prior Knowledge:
\begin{itemize}
\item Knowledge about Heterogenious computing and that memory needs to be accessible by the compute resource.
\end{itemize}
  What is introduced in this chapter:
\begin{itemize}
\item Various methods for giving the compute resource access to the memory it is working on
\item How memory in different physical spaces is allocated and deallocated
\item The next section only addresses how to move memory between different compute resources.
\item Movement can be asynchronous as long as it is available when it is used by the compute resource
\item Managing memory lifetime is seen as a separate problem.
\end{itemize}
}

\subsection{Zero-Copy Memory}
\gist{
  Prior Knowledge:
\begin{itemize}
\item that there is a way to move memory content between different physical memory locations.
\end{itemize}
  What is introduced in this chapter:
\begin{itemize}
\item values are read by another compute resource and sent directly via the interconnect to the consuming compute resource
\item the other way round for writing to memory.
\end{itemize}
}


\subsection{Explicit Memory Management}
\gist{
  Prior Knowledge:
\begin{itemize}
\item that there is a way to move memory content between different physical memory locations.
\end{itemize}
  What is introduced in this chapter:
\begin{itemize}
\item Explicit memory management means the user has to specify when to move which memory to what memory location.
\item Drivers adhere to the order memory movement and compute operations are given in
\end{itemize}
}

\subsection{Page Migration}
\gist{
  Prior Knowledge:
\begin{itemize}
\item Memory in computers is organized in memory pages and virtual addresses.
\item Which physical address a virtual address resides in is runtime dependent.
\end{itemize}
  What is introduced in this chapter:
\begin{itemize}
\item There exist multiple ways of resolving which physical location a page currently resides in.
\item Pages can be moved to another physical memory that is accessible by the requesting compute resource
\item This works via looking up the address in a address translation table and moving the page to a memory the compute resource can access
\item differentiate to unified memory here!
\item Do we need to also discuss that this can be just copying the data and not the page?
\item Pages stay where they are moved?
\item There is a unified address table that maps to all physical hardware memory locations.
\end{itemize}
}

\subsubsection{Address Translation in Hardware}
\gist{
  Prior Knowledge:
\begin{itemize}
\item Addresses need to be translated to pages and the pages need to be located.
\end{itemize}
  What is introduced in this chapter:
\begin{itemize}
\item Hardware address translation allows a compute resource to efficiently compute the page address of a memory address.
\item The driver can trigger a page migration if necessary.
\end{itemize}
}

\subsubsection{Address Translation in Software}
\gist{
  Prior Knowledge:
\begin{itemize}
\item Addresses need to be translated to pages and the pages need to be located.
\end{itemize}
  What is introduced in this chapter:
\begin{itemize}
\item Software address translation does not require specialized hardware and thus can be activated for older hardware by modifying driver and os implementation details.
\end{itemize}
}

\subsection{Unified Memory}
\gist{
  Prior Knowledge:
\begin{itemize}
\item Unified virtual address spaces.
\end{itemize}
  What is introduced in this chapter:
\begin{itemize}
\item Unified memory refers to hardware memory that can be directly accessed by all compute resources.
\item Also uses Unified virtual addresses.
\end{itemize}
}

\section{Kokkos' Memory Management Interface}



\section{GH200}
\import{./src_plots/}{vector_ping_pong_gh200.tex}

\section{MI300A}
\import{./src_plots/}{vector_ping_pong_mi300a.tex}

\end{document}
